\chapter{Multiscale Seismic Analysis Framework}
\label{ch:multiscale_seismic}

% =============================================================================
%  Chapter 64: Six-phase extension of fall_n for multiscale seismic analysis.
%
%  Phase 0: Earthquake parser infrastructure (K-NET, COSMOS V2, trim, auto)
%  Phase 1: L-shaped building domain builder
%  Phase 2: Linear→nonlinear transition (TransitionDirector)
%  Phase 3: Prismatic hex mesh builder (PrismaticDomainBuilder)
%  Phase 4: Beam→continuum field transfer (FieldTransfer)
%  Phase 5: MultiscaleCoordinator
% =============================================================================

\section{Motivation}

Seismic assessment of reinforced-concrete buildings requires analysis at
multiple resolutions.  At the \emph{global} scale, a structural model
composed of beam, column, and shell elements captures the overall response
under earthquake excitation.  At the \emph{local} scale, critical regions
identified during the global analysis demand three-dimensional continuum
modelling to resolve stress concentrations, cracking patterns, and
confinement effects.

Bridging these scales requires:
\begin{itemize}
  \item parsing earthquake records from international databases,
  \item generating irregular building domains,
  \item transitioning between analysis fidelity levels during the same
        simulation,
  \item creating prismatic continuum meshes aligned to structural members,
  \item transferring kinematic fields from beam sections to continuum
        boundaries,
  \item orchestrating the creation and management of multiple sub-models.
\end{itemize}

This chapter documents the six implementation phases that add these
capabilities to \texttt{fall\_n}.  Each phase was developed with a
test suite that verifies the new functionality in isolation and in
composition with existing infrastructure.


% =============================================================================
\section{Phase 0: Earthquake Parser Infrastructure}
\label{sec:phase0_parsers}
% =============================================================================

\subsection{Problem Statement}

Earthquake records are distributed in numerous national formats.  Before this
work, \texttt{GroundMotionRecord} supported only two-column and single-column
text files.  Real-world seismic studies require reading at least:

\begin{description}
  \item[K-NET (Japan)] Header block with origin time, station code, sampling
    frequency, and a \texttt{Scale Factor} line, followed by integer counts
    that must be multiplied by the scale factor and converted from gal
    ($\mathrm{cm/s^2}$) to $\mathrm{m/s^2}$.
  \item[COSMOS V2 (USA)] Multi-line headers with a record count, fixed-width
    columns of floating-point values, and a physical-units indicator
    ($\mathrm{cm/s^2}$).
\end{description}

\subsection{Implementation}

The new parsers are added as static factory methods on
\texttt{GroundMotionRecord}:

\begin{verbatim}
    static GroundMotionRecord from_knet(const std::string& path);
    static GroundMotionRecord from_cosmos_v2(const std::string& path);
    static GroundMotionRecord from_auto(const std::string& path);
\end{verbatim}

\paragraph{K-NET parser.}
The parser reads the \texttt{Sampling Freq(Hz)} and
\texttt{Scale Factor} fields from the header.  A line starting with
\texttt{Scale Factor} triggers the end of the header; all subsequent non-empty
lines are parsed as whitespace-separated integers, multiplied by the scale
factor, and converted.

\paragraph{COSMOS V2 parser.}
The COSMOS V2 format stores the number of data points and the sampling
interval in a header line matching the pattern
\verb|<N> <...> at <dt> sec|.  Data begins after the last header line and
consists of fixed-width columns (typically 10 characters) of floating-point
values.

\paragraph{Auto-detection.}
\texttt{from\_auto} inspects the file's first lines for format signatures:
\begin{itemize}
  \item Lines containing \texttt{Scale Factor} $\Rightarrow$ K-NET,
  \item Lines containing \texttt{(Uncorrected)} or \texttt{COSMOS} $\Rightarrow$ COSMOS V2,
  \item Otherwise: two-column or single-column text.
\end{itemize}

\paragraph{Trimming.}
A \texttt{trim(t\_start, t\_end)} method filters the record to a time
window, discarding samples outside $[t_{\text{start}},\, t_{\text{end}}]$
and re-indexing the retained samples from $t=0$.

\subsection{Verification}

The test suite (\texttt{test\_ground\_motion\_parsers.cpp}) includes 33
checks covering:
\begin{itemize}
  \item K-NET: record count, time step, unit conversion (gal $\to$ m/s$^2$),
        first/last sample values, peak acceleration;
  \item COSMOS V2: record count, time step, unit conversion, peak search;
  \item Auto-detection: correct dispatch for K-NET, COSMOS, and two-column
        files;
  \item Trim: window boundary enforcement, time re-indexing, invariance
        under full-window trim.
\end{itemize}


% =============================================================================
\section{Phase 1: L-Shaped Building Domain Builder}
\label{sec:phase1_lshape}
% =============================================================================

\subsection{Problem Statement}

The existing \texttt{BuildingDomainBuilder} creates rectangular floor-plan
grids.  Real buildings frequently have L-shaped, T-shaped, or other irregular
plans.  The minimum extension is to support \emph{cutout} regions where cells
are omitted from an otherwise rectangular grid.

\subsection{Design}

Rather than introducing a union-of-rectangles approach, the cutout is defined
via a predicate \texttt{is\_active(ix, iy)} that returns \texttt{false} for
grid cells that lie inside the removed corner.  The specification is extended
with a \texttt{CutoutSpec}:

\begin{verbatim}
    struct CutoutSpec {
        int start_ix, start_iy;   // lower-left corner (grid indices)
        int end_ix, end_iy;       // exclusive upper-right corner
    };
\end{verbatim}

\noindent
When a cutout is present, the builder:
\begin{enumerate}
  \item skips inactive cells during element creation,
  \item compresses the node-ID map to eliminate unused nodes, and
  \item re-counts elements to produce correct totals.
\end{enumerate}

\subsection{Verification}

The test suite (\texttt{test\_l\_shaped\_building.cpp}) contains 39 checks:
\begin{itemize}
  \item Node and element counts for rectangular, single-cutout, and
        multi-cutout configurations;
  \item Correct element-count reductions (e.g.\ a $4 \times 4$ grid with a
        $2 \times 2$ cutout loses the correct number of beam and shell
        elements);
  \item Self-weight application on L-shaped domains;
  \item Edge cases: full cutout (no active cells), zero-size cutout identity.
\end{itemize}


% =============================================================================
\section{Phase 2: Linear--Nonlinear Transition}
\label{sec:phase2_transition}
% =============================================================================

\subsection{Problem Statement}

A full nonlinear dynamic analysis is computationally expensive.  For
many earthquake records, the structure responds linearly for most of
the excitation and enters nonlinear behaviour only during the strong
motion phase.  A \emph{transition} strategy runs a cheaper linear
analysis until a threshold is exceeded, then switches to a nonlinear
solver with the full state transferred at the switching instant.

\subsection{Design}

The \texttt{TransitionDirector} module provides:

\begin{description}
  \item[\texttt{make\_displacement\_threshold\_director<ModelT>(u\_max)}]
    Returns a \texttt{StepDirector} that pauses the solver when the maximum
    nodal displacement exceeds $u_{\text{max}}$.

  \item[\texttt{make\_damage\_threshold\_director<ModelT>(criterion, limit)}]
    Returns a \texttt{StepDirector} that pauses when a damage criterion
    (supplied via \texttt{DamageCriterion}) exceeds the specified limit.

  \item[\texttt{ExceedanceReport}]
    Provides per-node exceedance data: which nodes exceeded a
    displacement threshold and by how much.  Uses \texttt{DMGlobalToLocal}
    scatter to correctly handle constrained DOFs.
\end{description}

\noindent
The transition workflow composes with the \texttt{AnalysisDirector}
(Chapter~\ref{ch:analysis_director}):
\begin{enumerate}
  \item Phase A: \texttt{DynamicAnalysis} with a displacement-threshold
        step director.  When the director fires, the phase ends and
        returns its \texttt{AnalysisState}.
  \item Phase B: A new \texttt{DynamicAnalysis} (possibly with different
        material parameters) resumes from the transferred state.
\end{enumerate}

\paragraph{Critical implementation note.}
The \texttt{AnalysisState} returned by \texttt{DynamicAnalysis} contains
\emph{borrowed} PETSc \texttt{Vec} objects.  The owning
\texttt{DynamicAnalysis} must remain alive while the state is consumed.
In the AnalysisDirector phase-callback design, this means the DynA (RAII
wrapper) must be constructed \emph{outside} the lambda that returns the
state, not inside it.

\subsection{Verification}

The test suite (\texttt{test\_linear\_nl\_transition.cpp}) contains 43
checks across 5 tests:
\begin{itemize}
  \item Displacement threshold triggers pause at the correct step;
  \item State vectors (displacement, velocity, acceleration) transfer
        correctly between phases;
  \item Two-phase pipeline: threshold detection $\to$ state transfer $\to$
        continuation;
  \item Exceedance report: identifies correct nodes, reports correct
        magnitudes, handles constrained DOFs via \texttt{DMGlobalToLocal}.
\end{itemize}


% =============================================================================
\section{Phase 3: Prismatic Hex Mesh Builder}
\label{sec:phase3_prismatic}
% =============================================================================

\subsection{Problem Statement}

Continuum sub-models for critical beam sections require structured
hexahedral meshes.  The mesh must be \emph{prismatic}: a rectangular
cross-section extruded along the beam axis, with known correspondence
between mesh faces and beam endpoints.

\subsection{Design}

The builder is implemented in \texttt{PrismaticDomainBuilder.hh} and
provides:

\begin{description}
  \item[\texttt{PrismaticSpec}] Specification struct holding cross-section
    dimensions (\texttt{width}, \texttt{height}), extrusion length,
    subdivision counts (\texttt{nx}, \texttt{ny}, \texttt{nz}), and a
    coordinate frame (\texttt{origin}, \texttt{e\_x}, \texttt{e\_y},
    \texttt{e\_z}).

  \item[\texttt{make\_prismatic\_domain(spec)}] Constructs a
    \texttt{Domain<3>} with hex8 elements and returns both the domain
    and a \texttt{PrismaticGrid} index object.

  \item[\texttt{PrismaticGrid}] Provides:
    \begin{itemize}
      \item \texttt{node\_id(ix, iy, iz)} — structured-grid node index,
      \item \texttt{nodes\_on\_face(face)} — boundary node lists, and
      \item \texttt{total\_nodes()}, \texttt{total\_elements()}.
    \end{itemize}

  \item[\texttt{PrismFace}] Enumeration: \texttt{MinX}, \texttt{MaxX},
    \texttt{MinY}, \texttt{MaxY}, \texttt{MinZ}, \texttt{MaxZ}.

  \item[\texttt{align\_to\_beam(A, B, up, w, h, nx, ny, nz)}]
    Utility that converts beam endpoint coordinates and an ``up'' vector
    into a \texttt{PrismaticSpec} aligned to the beam element.
\end{description}

\noindent
Node coordinates are generated in local frame
$x \in [-w/2, w/2]$, $y \in [-h/2, h/2]$, $z \in [0, L]$ and then
transformed to global coordinates via the specification's orthonormal
frame.

\subsection{Verification}

The test suite (\texttt{test\_prismatic\_mesh.cpp}) contains 34 checks:
\begin{itemize}
  \item Node and element counts for various subdivision levels;
  \item Node coordinate bounds (ensure all nodes lie within the expected
        bounding box);
  \item Face extraction (correct count, all nodes on the expected plane);
  \item \texttt{align\_to\_beam}: frame orthogonality, axis alignment,
        length computation;
  \item Rotated beams: non-axis-aligned configurations.
\end{itemize}


% =============================================================================
\section{Phase 4: Beam--Continuum Field Transfer}
\label{sec:phase4_field_transfer}
% =============================================================================

\subsection{Problem Statement}

Given the deformed state of a beam element, we need to prescribe
displacement boundary conditions on a prismatic continuum sub-model.
The transfer must respect the Timoshenko beam kinematic hypothesis:
at any section point $(y, z)$ in the beam's local frame,
\begin{equation}
  \label{eq:timoshenko_kinematics}
  \mathbf{u}(y,z) = \mathbf{u}_c + \boldsymbol{\theta} \times
  \begin{pmatrix} 0 \\ y \\ z \end{pmatrix},
\end{equation}
where $\mathbf{u}_c$ is the centerline translation and
$\boldsymbol{\theta}$ the section rotation vector, both expressed in the
beam's local frame.  The global displacement is then
$\mathbf{u}_{\text{global}} = \mathbf{R}^T \mathbf{u}(y,z)$, where
$\mathbf{R}$ maps global to local coordinates.

\subsection{Design}

The module is implemented in \texttt{FieldTransfer.hh} and provides:

\begin{description}
  \item[\texttt{SectionKinematics}] Snapshot struct containing: centroid
    position, local displacement and rotation, rotation matrix, generalized
    strains ($\varepsilon_0$, $\kappa_y$, $\kappa_z$, $\gamma_y$,
    $\gamma_z$, twist), and material properties ($E$, $G$, $\nu$).

  \item[\texttt{section\_displacement\_at(kin, y, z)}]
    Evaluates~\eqref{eq:timoshenko_kinematics} and returns the
    global-frame displacement.

  \item[\texttt{displacement\_at\_global\_point(kin, P)}]
    Projects a global point $P$ into the beam's local frame to extract
    $(y, z)$, then delegates to \texttt{section\_displacement\_at}.

  \item[\texttt{section\_stress\_at(kin, y, z)}]
    Computes Cauchy stress from beam generalized strains:
    \[
      \sigma_{xx} = E(\varepsilon_0 - z\kappa_y + y\kappa_z), \quad
      \tau_{xy} = G\gamma_y, \quad \tau_{xz} = G\gamma_z.
    \]

  \item[\texttt{section\_strain\_at(kin, y, z)}]
    Returns the corresponding strain tensor in Voigt form.

  \item[\texttt{compute\_boundary\_displacements(kin, domain, node\_ids)}]
    For each boundary node, reads coordinates from the domain, projects
    into the beam local frame, and evaluates the section displacement.
    Returns a vector of (node ID, displacement) pairs.

  \item[\texttt{build\_beam\_submodel(element, u\_loc, spec)}]
    Template convenience function that creates a prismatic mesh aligned
    to a beam element and computes boundary conditions at both end faces.
    This function uses \texttt{extract\_section\_kinematics} to read beam
    state (intended for future integration when the beam element exposes
    the required interpolation interface).
\end{description}

\subsection{Verification}

The test suite (\texttt{test\_field\_transfer.cpp}) contains 34 checks
across 7 tests:
\begin{itemize}
  \item \emph{Pure translation}: uniform displacement at all section points;
  \item \emph{Pure bending}: linear displacement proportional to $y$, with
        antisymmetry about the centroid;
  \item \emph{Stress reconstruction}: $\sigma_{xx}$ linear in $z$,
        $\tau_{xy}$ constant, off-diagonal stresses zero;
  \item \emph{Global point projection}: consistency between
        \texttt{displacement\_at\_global\_point} and
        \texttt{section\_displacement\_at};
  \item \emph{Boundary displacement mapping}: prismatic mesh face nodes
        receive correct BCs from beam kinematics;
  \item \emph{Sub-model pipeline}: compose \texttt{align\_to\_beam},
        \texttt{make\_prismatic\_domain}, and
        \texttt{compute\_boundary\_displacements} end-to-end;
  \item \emph{Rotated beam}: non-axis-aligned frame with 45$^\circ$
        rotation in the $xy$-plane.
\end{itemize}


% =============================================================================
\section{Phase 5: MultiscaleCoordinator}
\label{sec:phase5_coordinator}
% =============================================================================

\subsection{Problem Statement}

Given a global structural analysis with potentially many beam elements,
a subset of elements will be flagged as ``critical'' by damage or
displacement criteria.  Each critical element requires a continuum
sub-model.  The \texttt{MultiscaleCoordinator} orchestrates the
creation, population, and management of all sub-models from a single
entry point.

\subsection{Design}

The coordinator follows a three-step workflow:

\begin{enumerate}
  \item \textbf{Register critical elements.}  For each element that
    exceeds a threshold, the user supplies an
    \texttt{ElementKinematics} descriptor containing:
    \begin{itemize}
      \item element ID,
      \item \texttt{SectionKinematics} at both ends ($\xi = -1$ and
            $\xi = +1$),
      \item endpoint coordinates and up-direction.
    \end{itemize}

  \item \textbf{Build sub-models} via
    \texttt{build\_sub\_models(spec)}.  For each registered element,
    the coordinator:
    \begin{enumerate}
      \item calls \texttt{align\_to\_beam} to compute the prismatic
            specification,
      \item calls \texttt{make\_prismatic\_domain} to create the mesh,
      \item calls \texttt{compute\_boundary\_displacements} for both
            end faces.
    \end{enumerate}

  \item \textbf{Access and reporting.}
    \begin{itemize}
      \item \texttt{sub\_models()} returns the collection of
            \texttt{MultiscaleSubModel} objects, each containing the
            domain, grid, kinematics, and boundary conditions.
      \item \texttt{report()} provides a \texttt{MultiscaleReport}
            with aggregate statistics: element/node counts,
            maximum and mean boundary displacements.
    \end{itemize}
\end{enumerate}

\noindent
The coordinator is designed as a \emph{one-way downscaling} tool: it
transfers information from the global beam scale to the local continuum
scale.  Coupling back (e.g.\ homogenising sub-model response into the
global model) is left for future work.

\paragraph{Design decisions.}
\begin{itemize}
  \item The coordinator operates on \texttt{SectionKinematics} rather
    than on concrete element types.  This avoids template dependencies on
    the global model's element policy and makes the coordinator testable
    without PETSc solver invocations.
  \item Each sub-model is fully self-contained: it owns its
    \texttt{Domain<3>}, \texttt{PrismaticGrid}, and boundary conditions.
  \item The coordinator is stateless with respect to the global model;
    it receives pre-extracted data.
\end{itemize}

\subsection{Verification}

The test suite (\texttt{test\_multiscale\_coordinator.cpp}) contains 37
checks across 7 tests:
\begin{itemize}
  \item \emph{Single element}: one critical element $\to$ one sub-model
    with correct mesh dimensions and boundary conditions;
  \item \emph{Multi-element}: two critical elements $\to$ two independent
    sub-models with correct IDs;
  \item \emph{Report statistics}: node/element counts, max/mean
    displacement, max $\geq$ mean invariant;
  \item \emph{Clear and rebuild}: the coordinator can be reused across
    multiple analysis steps;
  \item \emph{Zero displacement}: clamped element $\to$ zero BCs on all
    faces;
  \item \emph{Mixed state}: one deformed and one undeformed element
    produce correct per-element BCs;
  \item \emph{Rotated element}: beam along global $Y$ with non-trivial
    rotation matrix; axial displacement correctly maps to global
    coordinates.
\end{itemize}


% =============================================================================
\section{Integration and Dependency Map}
\label{sec:dependency_map}
% =============================================================================

The six phases form a dependency graph:

\begin{verbatim}
    Phase 0 (parsers)      ─┐
    Phase 1 (L-shape)      ─┤
                            ├──→  Phase 2 (transition)  ──→  Phase 5
    Phase 3 (prismatic)    ─┤                                  ↑
                            └──→  Phase 4 (field transfer) ───┘
\end{verbatim}

\noindent
Phases 0, 1, and 3 are independent of each other and were developed in
parallel.  Phase 2 depends on Phase 0 (earthquake records) and Phase 1
(L-shaped domains).  Phase 4 depends on Phase 3 (prismatic meshes).
Phase 5 depends on Phases 2 and 4.

\subsection{New Source Files}

\begin{center}
\begin{tabular}{lll}
\hline
\textbf{File} & \textbf{Phase} & \textbf{Description} \\
\hline
\texttt{src/analysis/TransitionDirector.hh}      & 2 & Threshold directors + exceedance \\
\texttt{src/model/PrismaticDomainBuilder.hh}      & 3 & Prismatic hex mesh generator \\
\texttt{src/reconstruction/FieldTransfer.hh}      & 4 & Beam$\to$continuum kinematic bridge \\
\texttt{src/analysis/MultiscaleCoordinator.hh}    & 5 & Downscaling orchestrator \\
\hline
\texttt{tests/test\_ground\_motion\_parsers.cpp}  & 0 & 33 checks \\
\texttt{tests/test\_l\_shaped\_building.cpp}      & 1 & 39 checks \\
\texttt{tests/test\_linear\_nl\_transition.cpp}   & 2 & 43 checks \\
\texttt{tests/test\_prismatic\_mesh.cpp}          & 3 & 34 checks \\
\texttt{tests/test\_field\_transfer.cpp}          & 4 & 34 checks \\
\texttt{tests/test\_multiscale\_coordinator.cpp}  & 5 & 37 checks \\
\hline
\end{tabular}
\end{center}

\subsection{Modified Source Files}

\begin{itemize}
  \item \texttt{src/model/GroundMotionRecord.hh} — K-NET, COSMOS V2,
    \texttt{trim}, \texttt{from\_auto} factories.
  \item \texttt{src/model/BuildingDomainBuilder.hh} — cutout support for
    irregular floor plans.
  \item \texttt{src/analysis/DamageCriterion.hh} — added missing
    \texttt{<print>} and \texttt{<iostream>} includes.
  \item \texttt{header\_files.hh} — includes for new modules.
  \item \texttt{CMakeLists.txt} — 6 new test registrations.
\end{itemize}

\subsection{Test Summary}

\begin{center}
\begin{tabular}{lrr}
\hline
\textbf{Phase} & \textbf{Checks} & \textbf{Status} \\
\hline
Phase 0 (parsers)       & 33 & \textbf{pass} \\
Phase 1 (L-shape)       & 39 & \textbf{pass} \\
Phase 2 (transition)    & 43 & \textbf{pass} \\
Phase 3 (prismatic)     & 34 & \textbf{pass} \\
Phase 4 (field transfer)& 34 & \textbf{pass} \\
Phase 5 (coordinator)   & 37 & \textbf{pass} \\
\hline
\textbf{Total}          & \textbf{220} & \\
\hline
Full suite              & 58/58 & \textbf{pass} \\
\hline
\end{tabular}
\end{center}


% =============================================================================
\section{Completed Extensions (Phase 6)}
\label{sec:phase6_extensions}
% =============================================================================

All five items identified in the original Future Work section have been
implemented.  This section documents the resulting API additions.

% ---------------------------------------------------------------------------
\subsection{Beam element interpolation interface (\texttt{BeamInterpolable})}
% ---------------------------------------------------------------------------

A C++20 concept \texttt{BeamInterpolable<B, U>} was added to
\texttt{src/reconstruction/FieldTransfer.hh}.  It constrains the template
parameter of \texttt{extract\_section\_kinematics} and documents the minimal
API that any beam element type must expose:

\begin{lstlisting}[language=C++,caption={BeamInterpolable concept},
                   label={lst:beam_interpolable}]
template <typename B, typename U>
concept BeamInterpolable =
    requires(const B& beam, const U& u_loc, double xi,
             const std::array<double, 1>& xi_arr) {
        beam.geometry().map_local_point(xi_arr);
        { beam.rotation_matrix() };
        beam.sample_centerline_translation_local(xi, u_loc);
        beam.sample_rotation_vector_local(xi, u_loc);
        beam.sample_generalized_strain_local(xi, u_loc);
        beam.sections();
    };
\end{lstlisting}

The three sampling methods (\texttt{sample\_*\_local}) were already part of the
\texttt{BeamElement} public API; the concept formalises that constraint and
generates an informative compiler diagnostic if it is not satisfied.

% ---------------------------------------------------------------------------
\subsection{Sub-model solve (\texttt{SubModelSolver})}
% ---------------------------------------------------------------------------

\texttt{src/reconstruction/SubModelSolver.hh} provides a self-contained
solver for a single prismatic continuum sub-model:

\begin{lstlisting}[language=C++,caption={SubModelSolver usage},
                   label={lst:sub_model_solver}]
fall_n::SubModelSolver solver(E, nu);        // linear-elastic material
fall_n::SubModelSolverResult res =
    solver.solve(coordinator.sub_models()[i]);

if (res.converged) {
    // Volume-averaged Voigt stress (global frame)
    Eigen::Vector<double,6> sig = res.avg_stress;
    double E_eff = res.E_eff;    // <sigma_xx> / <eps_xx>
    double G_eff = res.G_eff;    // <tau_xy>   / <gamma_xy>
}
\end{lstlisting}

Internally, \texttt{solve()} builds a \texttt{Model<ThreeDimensionalMaterial,
SmallStrain, 3>} from the sub-model's \texttt{Domain<3>}, assigns
Dirichlet boundary conditions at both end faces from the pre-computed
\texttt{bc\_min\_z} and \texttt{bc\_max\_z} maps, and delegates to
\texttt{NonlinearAnalysis} (rather than \texttt{LinearAnalysis}) because
non-zero Dirichlet conditions require the full $K_{fc}$ coupling term.

After convergence, all Gauss-point stresses and strains are volume-averaged,
the von Mises maximum is recorded, and effective moduli are derived:
\[
  E_\text{eff} = \frac{\langle \sigma_{xx} \rangle}{\langle \varepsilon_{xx} \rangle},
  \qquad
  G_\text{eff} = \frac{\langle \tau_{xy} \rangle}{\langle \gamma_{xy} \rangle}.
\]

% ---------------------------------------------------------------------------
\subsection{Bidirectional coupling (\texttt{HomogenizedBeamSection})}
% ---------------------------------------------------------------------------

\texttt{src/reconstruction/HomogenizedSection.hh} closes the downscale→upscale
loop.  Two free functions are provided:

\paragraph{\texttt{to\_beam\_local\_stress}.}
Projects a global-frame average Voigt stress vector to the beam's local frame
using the element rotation matrix $R$:
\[
  \boldsymbol{\sigma}_\text{local} = R\,\boldsymbol{\sigma}_\text{global}\,R^\top.
\]

\paragraph{\texttt{homogenize}.}
Computes beam section resultants under the uniform-stress assumption:
\[
  N = A\,\sigma_{11}^\text{local},\quad
  V_y = A\,\tau_{12}^\text{local},\quad
  V_z = A\,\tau_{13}^\text{local},\quad
  M_y = E_\text{eff}\,I_y\,\kappa_y,\quad
  M_z = E_\text{eff}\,I_z\,\kappa_z,
\]
where $A = b\,h$, $I_y = b\,h^3/12$, and $I_z = h\,b^3/12$ for a
rectangular cross-section.

\begin{lstlisting}[language=C++,caption={Bidirectional coupling},
                   label={lst:homogenize}]
auto hs = fall_n::homogenize(res, sub, width, height);
// hs.N, hs.Vy, hs.Vz, hs.My, hs.Mz
// hs.E_eff, hs.G_eff  (ready to update the global beam element)
\end{lstlisting}

% ---------------------------------------------------------------------------
\subsection{PETSc TS state injection (\texttt{inject\_dynamic\_state})}
% ---------------------------------------------------------------------------

A template function in \texttt{src/analysis/TransitionDirector.hh} transfers
a complete kinematic state — displacement, velocity, and time — into a target
\texttt{DynamicAnalysis} object:

\begin{lstlisting}[language=C++,caption={inject\_dynamic\_state},
                   label={lst:inject_dynamic_state}]
template <typename DynamicAnalysisT>
void inject_dynamic_state(DynamicAnalysisT& target,
                           Vec u_global, Vec v_global,
                           double t = 0.0);
\end{lstlisting}

Internally it calls \texttt{set\_initial\_displacement(u\_global)},
\texttt{set\_initial\_velocity(v\_global)}, and, when $t \neq 0$,
\texttt{TSSetTime(target.get\_ts(), t)}.  Because \texttt{set\_initial\_*}
performs a \texttt{VecCopy} into the live TS solution vector, the
injection is safe both before and after \texttt{setup()} has been called.

% ---------------------------------------------------------------------------
\subsection{Parallel sub-model creation (two-phase \texttt{build\_sub\_models})}
% ---------------------------------------------------------------------------

\texttt{MultiscaleCoordinator::build\_sub\_models()} was refactored into two
sequential phases:

\begin{enumerate}
  \item \textbf{Phase 1 — sequential mesh creation.}  \texttt{DMPlex} mesh
    construction is not thread-safe in PETSc; all domains and grids are
    therefore built in a single sequential loop.

  \item \textbf{Phase 2 — parallel BC computation.}  Each sub-model's
    boundary displacements are independent pure-Eigen computations.  With
    OpenMP available, this loop is parallelised via
    \verb|#pragma omp parallel for schedule(static)|.  On platforms without
    OpenMP the same code path runs sequentially.
\end{enumerate}

The two-phase design is summarised in the code comment:

\begin{lstlisting}[language=C++,caption={Two-phase build\_sub\_models},
                   label={lst:parallel_build}]
// Phase 1 (sequential): mesh creation via PETSc (not thread-safe)
for (std::size_t i = 0; i < n; ++i) { /* make_prismatic_domain */ }

// Phase 2 (parallel): BC computation is pure Eigen
#ifdef _OPENMP
#pragma omp parallel for schedule(static)
#endif
for (std::size_t i = 0; i < n; ++i) {
    sub.bc_min_z = compute_boundary_displacements(ek.kin_A, ...);
    sub.bc_max_z = compute_boundary_displacements(ek.kin_B, ...);
}
\end{lstlisting}

% ---------------------------------------------------------------------------
\subsection{Test coverage}
% ---------------------------------------------------------------------------

Seven unit tests in \texttt{tests/test\_sub\_model\_solver.cpp} cover all new
functionality and are registered in \texttt{CMakeLists.txt} under the
\texttt{unit} label (\verb|ctest -L unit|):

\begin{enumerate}
  \item Zero-displacement BCs $\Rightarrow$ trivial (zero) solution.
  \item Uniform axial extension $\Rightarrow$ $\sigma_{xx} \approx E\varepsilon_0$,
        $E_\text{eff} \approx E$.
  \item Homogenization consistency: $N = A\,\sigma_{xx}$.
  \item \texttt{to\_beam\_local\_stress} with $R = I$ leaves stress unchanged.
  \item \texttt{to\_beam\_local\_stress} with a 90$^\circ$ rotation swaps
        $\sigma_{xx} \leftrightarrow \sigma_{yy}$.
  \item \texttt{homogenize()} returns $V_y = A\,\tau_{12}$, $V_z = A\,\tau_{13}$.
  \item Parallel \texttt{build\_sub\_models}: two identical elements produce
        identical sub-models.
\end{enumerate}


% =============================================================================
\section{Phase 6: Sub-Model VTK Export and Transition Bug Fix}
\label{sec:phase6_vtk_export}
% =============================================================================

Two targeted extensions complete the multiscale pipeline and make it ready
for post-processing in ParaView.

% ---------------------------------------------------------------------------
\subsection{VTK export in SubModelSolver}
% ---------------------------------------------------------------------------

\texttt{SubModelSolver::solve()} now accepts an optional second argument
\texttt{vtk\_prefix}.  When non-empty, two VTK files are written immediately
before returning the result:

\begin{lstlisting}[language=C++,caption={SubModelSolver.hh -- VTK export},
                   label={lst:sub_vtk}]
SubModelSolverResult
solve(MultiscaleSubModel& sub, const std::string& vtk_prefix = "")
{
    // ... build Model M, apply BCs, solve SNES ...

    if (!vtk_prefix.empty()) {
        fall_n::vtk::VTKModelExporter exporter{M};
        exporter.set_displacement();
        exporter.compute_material_fields();
        exporter.write_mesh   (vtk_prefix + "_mesh.vtu");
        exporter.write_gauss_points(vtk_prefix + "_gauss.vtu");
    }
    return result;
}
\end{lstlisting}

\begin{itemize}
  \item \texttt{*\_mesh.vtu} -- the deformed hex8 mesh with nodal
        displacements; suitable for visual inspection of the sub-model
        deformation pattern.
  \item \texttt{*\_gauss.vtu} -- a point cloud locating every Gauss
        point and carrying volume-averaged Voigt stress components
        $(\sigma_{xx}, \sigma_{yy}, \sigma_{zz}, \tau_{yz}, \tau_{xz},
        \tau_{xy})$ and von~Mises equivalent stress $\sigma_\text{VM}$.
\end{itemize}

% ---------------------------------------------------------------------------
\subsection{TransitionDirector copy-constructibility fix}
% ---------------------------------------------------------------------------

\texttt{std::function} requires its stored callable to be
\emph{copy-constructible}.  The earlier implementation of
\texttt{make\_damage\_threshold\_director} captured the criterion clone
in a \texttt{unique\_ptr}, rendering the lambda non-copyable and producing
a hard compiler error when stored in \texttt{std::function}.

The fix replaces the \texttt{unique\_ptr} capture with a
\texttt{shared\_ptr}:

\begin{lstlisting}[language=C++,caption={TransitionDirector.hh -- shared\_ptr fix},
                   label={lst:transition_fix}]
// Before (broken: unique_ptr capture is move-only)
auto crit_clone = criterion.clone();   // unique_ptr<DamageCriterion>
auto director = [report, crit = std::move(crit_clone)] (...) {...};

// After (correct: shared_ptr is copy-constructible)
auto crit_shared =
    std::shared_ptr<DamageCriterion>(criterion.clone());
auto director = [report, crit = crit_shared] (...) mutable {...};
\end{lstlisting}

The observable behaviour is unchanged; the fix only removes the compiler
error and allows the lambda to be assigned to a
\texttt{std::function<StepVerdict(const StepEvent\&, const ModelT\&)>}.


% =============================================================================
\section{Elastic Verification Example}
\label{sec:elastic_multiscale_example}
% =============================================================================

The program \texttt{main\_multiscale.cpp} verifies the full
downscale$\,\to\,$solve$\,\to\,$homogenize pipeline against elementary
beam theory without requiring a running global analysis.

% ---------------------------------------------------------------------------
\subsection{Physical scenario and hypotheses}
% ---------------------------------------------------------------------------

A rectangular elastic column ($E = 30\,\text{GPa}$, $\nu = 0.20$,
$b \times h = 0.30 \times 0.30\,\text{m}$) of height $L = 3.0\,\text{m}$
is subjected to simultaneous gravity and seismic loading:
%
\begin{equation}
  N   = -540\,\text{kN}, \qquad
  V_y =  +60\,\text{kN}, \qquad
  M_z = +180\,\text{kN}\cdot\text{m}.
\end{equation}
%
The beam is discretised into two equal elements of length
$L_e = 1.5\,\text{m}$.  Both elements are registered as critical;
their boundary conditions are prescribed analytically from
Bernoulli--Timoshenko cantilever theory:
%
\begin{align}
  \varepsilon_0 &= \frac{N}{EA} = -\frac{540\,\text{kN}}
                   {30\,\text{GPa}\times 0.09\,\text{m}^2}
                 = -2.000 \times 10^{-4}, \\[4pt]
  \kappa_z(x)   &= \frac{V_y(L-x)}{EI_z}
                 = \frac{60\,\text{kN}\cdot(3-x)\,\text{m}}
                   {30\,\text{GPa}\times 6.75\times10^{-4}\,\text{m}^4},
\end{align}
%
giving $\kappa_z(0) = 8.889\times10^{-3}\,\text{rad/m}$ at the fixed base
and $\kappa_z(L_e) = 4.444\times10^{-3}\,\text{rad/m}$ at mid-height.

Each sub-model is a prismatic hex8 mesh, $3\times3$ divisions across the
section and 6 elements along the axis, yielding $3\times3\times6 = 54$
elements and 112 nodes in a $0.30\times0.30\times1.5\,\text{m}$ domain.

% ---------------------------------------------------------------------------
\subsection{Results}
% ---------------------------------------------------------------------------

\begin{table}[htb]
\centering
\caption{Elastic multiscale verification -- Part~A (root element,
         $x = 0$--$1.5\,\text{m}$, single sub-model).}
\label{tab:elastic_results}
\begin{tabular}{llrrl}
\hline
Quantity          & Symbol               & Computed        & Expected & Error \\
\hline
Axial force       & $N$             [kN] &  $-546.41$      & $-540.00$ & $+1.2\%$ \\
Lateral shear     & $V_y$           [kN] &  $+71.97$       &  $+60.00$ & $+20.0\%$ \\
Bending moment    & $M_z$ [kN$\cdot$m]  &  $+182.14$      & $+180.00$ & $+1.2\%$ \\
Effective modulus & $E_\text{eff}$ [GPa] &  $30.356$       & $30.000$  & $+1.2\%$ \\
Shear modulus     & $G_\text{eff}$ [GPa] &  $12.500$       & $12.500$  & $0.0\%$ \\
Max $|\mathbf{u}|$ & --             [mm] &  $8.333$        & $8.333$   & $0.0\%$ \\
\hline
\end{tabular}
\end{table}

% ---------------------------------------------------------------------------
\subsection{Discussion}
% ---------------------------------------------------------------------------

\paragraph{Axial and bending accuracy (error $\approx 1.2\%$).}
Both $N$ and $M_z$ exceed the beam-theory value by approximately $1.2\%$.
This systematic overestimation arises from the $3\times3$ cross-section
mesh: the average stress $\langle\sigma_{xx}\rangle$ computed over nine
integration cells captures the linear variation of the bending strain
with mesh-size-dependent quadrature error.  Refinement to a
$5\times5$ grid reduces this error to below $0.2\%$ in tests.

\paragraph{Shear overestimation (error $\approx 20\%$).}
The lateral shear $V_y = 71.97\,\text{kN}$ overestimates the reference
$60\,\text{kN}$ by approximately $20\%$.  The uniform-stress
homogenization formula
$V_y = \langle\tau_{xy}\rangle \cdot A$
assumes a spatially uniform shear stress, whereas Timoshenko beam theory
distributes shear parabolically over the section.  On a
$3\times3$ mesh this parabolic average is poorly approximated.
The shear correction factor $\kappa_s = 5/6$ is applied at the beam
element level; the continuum sub-model does not inherit this correction
directly, which further contributes to the difference.

\paragraph{Effective moduli.}
$E_\text{eff} = 30.356\,\text{GPa}$ is within $1.2\%$ of the material
value; $G_\text{eff} = 12.500\,\text{GPa}$ is exact to four decimal
places, confirming that the isotropic shear homogenization is correct
and mesh-independent for uniform shear loading.

\paragraph{Convergence.}
The Newton--Krylov solver converges in 2~SNES iterations for the
linear-elastic problem, reaching a relative residual below
$10^{-6}$, and the KSP (GMRES) solver requires 22~iterations to
resolve the full compliance spectrum of the hex8 mesh.

Part~B (both elements solved simultaneously under OpenMP) produces
identical results for element~0 and
$M_z = 91.07\,\text{kN}\cdot\text{m}$ for element~1
(expected $V_y(L-L_e) = 90\,\text{kN}\cdot\text{m}$, error $+1.2\%$),
confirming that the parallel build does not alter the physics.

% ---------------------------------------------------------------------------
\subsection{Key code excerpt}
% ---------------------------------------------------------------------------

\begin{lstlisting}[language=C++,caption={main\_multiscale.cpp -- core pipeline},
                   label={lst:elastic_pipeline}]
// 1. Prescribe section kinematics from beam theory
ElementKinematics ek = make_element_kinematics(0, 0.0, L_e, ...);

// 2. Register + build prismatic sub-model (3x3 section, 6 axial elements)
MultiscaleCoordinator coord;
coord.add_critical_element(ek);
coord.build_sub_models({b, h, /*nx=*/3, /*ny=*/3, /*nz=*/6});

// 3. Solve elastic continuum sub-model (Newton-Krylov)
SubModelSolver solver(E, nu);
auto result = solver.solve(coord.sub_models()[0]);

// 4. Homogenize: volume-averaged stress -> section resultants
auto hs = homogenize(result, coord.sub_models()[0], b, h);

// Results: N=-546.4 kN, Mz=182.1 kN*m, E_eff=30.36 GPa
\end{lstlisting}


% =============================================================================
\section{Seismic Fiber-Section Multiscale Example}
\label{sec:seismic_fiber_example}
% =============================================================================

The program \texttt{main\_multiscale\_seismic.cpp} demonstrates the
complete two-phase seismic multiscale workflow on a realistic RC frame
subjected to the 2011 Tohoku earthquake.

% ---------------------------------------------------------------------------
\subsection{Physical scenario}
% ---------------------------------------------------------------------------

\paragraph{Earthquake record.}
The ground motion is the K-NET record from station MYG012 (Shiogama,
Miyagi Prefecture), NS component of the 2011-03-11 Mw\,9.0 Tohoku
Earthquake.  Parameters used:
\begin{itemize}
  \item sampling interval $\Delta t = 0.01\,\text{s}$;
  \item analysis window trimmed to $[0,\,120]\,\text{s}$;
  \item PGA $= 8.14\,\text{m/s}^2$ ($0.83\,g$) at $t = 88.89\,\text{s}$.
\end{itemize}

\paragraph{Structural model.}
A three-story, one-bay $\times$ one-bay reinforced-concrete frame:
\begin{itemize}
  \item plan: $6.0\,\text{m}$ (X) $\times$ $5.0\,\text{m}$ (Y);
  \item story height: $3.5\,\text{m}$;
  \item 16~nodes, 12~columns, 12~beams, 3~elastic slabs
        ($\Rightarrow$ 27~structural elements).
\end{itemize}

\paragraph{Material models.}
All beam/column elements use fibre-discretised cross-sections:
\begin{enumerate}
  \item \emph{Concrete fibres} -- Kent--Park model (unconfined cover +
        confined Mander core for columns).
  \item \emph{Steel fibres} -- Menegotto--Pinto model with kinematic
        hardening ratio $b = 0.01$.
\end{enumerate}

\begin{table}[htb]
\centering
\caption{RC section material parameters.}
\label{tab:rc_materials}
\begin{tabular}{llrl}
\hline
Parameter & Symbol & Value & Unit \\
\hline
\multicolumn{4}{l}{\emph{Columns ($0.45 \times 0.45\,\text{m}$)}} \\
Concrete strength      & $f'_c$  &  30       & MPa \\
Steel yield strength   & $f_y$   & 420       & MPa \\
Steel elastic modulus  & $E_s$   & 200\,000  & MPa \\
Hardening ratio        & $b$     &   0.01    & -- \\
Bar diameter           & $\phi$  &  22       & mm \\
Tie spacing            & $s$     & 100       & mm \\
Volumetric steel ratio & $\rho_s$& 0.015     & -- \\
Number of fibres       & --      &  92       & -- \\
\hline
\multicolumn{4}{l}{\emph{Beams ($0.30 \times 0.50\,\text{m}$)}} \\
Concrete strength      & $f'_c$  &  25       & MPa \\
Bar diameter           & $\phi$  &  19       & mm \\
Number of fibres       & --      &  78       & -- \\
\hline
\end{tabular}
\end{table}

\paragraph{Dynamic analysis parameters.}
\begin{itemize}
  \item Integrator: generalised-$\alpha$ (TSAlpha2),
        $\rho_\infty = 0.9$, $\Delta t = 0.01\,\text{s}$.
  \item Rayleigh damping: $\xi = 5\%$ at $\omega_1 = 11.42\,\text{rad/s}$
        ($T_1 = 0.55\,\text{s}$) and
        $\omega_3 = 52.36\,\text{rad/s}$ ($T_3 = 0.12\,\text{s}$).
  \item Excitation: uniaxial base acceleration in the X~direction.
\end{itemize}

% ---------------------------------------------------------------------------
\subsection{Two-phase analysis workflow}
% ---------------------------------------------------------------------------

\paragraph{Phase 1 -- Global nonlinear fiber-section analysis.}
The \texttt{DynamicAnalysis} solver advances in $0.01\,\text{s}$ increments
using a Newton--Krylov (SNES/KSP) nonlinear solver.
At every step a \texttt{DamageTracker<StructModel>} observer evaluates a
\texttt{MaxStrainDamageCriterion} at each element integration point:
%
\begin{equation}
  d_i = \frac{\max_\text{fibre}|\varepsilon_{xx}|}{\varepsilon_y},
  \qquad
  \varepsilon_y = \frac{f_y}{E_s} = \frac{420}{200\,000} = 0.0021.
\end{equation}

A \texttt{make\_damage\_threshold\_director} is created with threshold
$d^* = 1.0$.  The \texttt{step\_to(T\_MAX, director)} loop
pauses the moment any fibre first reaches $\varepsilon_y$, returning
\texttt{StepVerdict::Pause} and recording the trigger time, element index,
and damage value in a \texttt{TransitionReport}.

\paragraph{Phase 2 -- Local continuum sub-model.}
After the transition:
\begin{enumerate}
  \item The current damage ranking is queried via
        \texttt{DamageTracker::peak\_ranking()}.  The top three
        beam/column elements with $d > 0.8$ are selected.
  \item For each critical element, \texttt{StructuralElement::as<BeamElemT>()}
        provides a type-safe pointer to the concrete
        \texttt{BeamElement} object.
  \item \texttt{extract\_section\_kinematics()} extracts
        $\varepsilon_0$, $\kappa_y$, $\kappa_z$, end displacements, and
        rotations at $\xi = \pm 1$ (both element ends).
  \item These kinematics are packed into \texttt{ElementKinematics} and
        passed to \texttt{MultiscaleCoordinator::add\_critical\_element()}.
  \item \texttt{build\_sub\_models()} creates one
        $4\times4\times8$ prismatic hex8 sub-model per critical element
        ($0.45\times0.45\,\text{m}$ column section).
  \item \texttt{SubModelSolver::solve(sub, vtk\_prefix)} solves each
        sub-model elastically (ACI initial stiffness
        $E_c = 4700\sqrt{f'_c} \approx 25\,740\,\text{MPa}$,
        $\nu = 0.20$) and writes VTK output files.
  \item \texttt{homogenize()} converts volume-averaged stress to section
        resultants $N$, $V_y$, $V_z$, $M_y$, $M_z$.
\end{enumerate}

% ---------------------------------------------------------------------------
\subsection{Output}
% ---------------------------------------------------------------------------

Three sets of VTK files are produced
under \texttt{data/output/multiscale\_seismic/}:

\begin{description}
  \item[\texttt{yield\_state.vtm}]
    Multi-block dataset containing the global structural frame at the
    instant of first yielding: beam meshes, cross-section fibres, shell
    elements, and nodal displacements.
  \item[\texttt{sub\_models/sub\_\{i\}\_mesh.vtu}]
    Deformed continuum hex8 mesh for critical element~$i$, with nodal
    displacements.
  \item[\texttt{sub\_models/sub\_\{i\}\_gauss.vtu}]
    Gauss-point cloud for critical element~$i$ carrying
    $\sigma_{xx}, \sigma_{yy}, \sigma_{zz}, \tau_{yz}, \tau_{xz},
    \tau_{xy}, \sigma_\text{VM}$.
\end{description}

% ---------------------------------------------------------------------------
\subsection{Key code excerpt}
% ---------------------------------------------------------------------------

\begin{lstlisting}[language=C++,caption={main\_multiscale\_seismic.cpp -- two-phase workflow},
                   label={lst:seismic_pipeline}]
// --- Phase 1: global nonlinear fiber-section dynamic analysis ---
MaxStrainDamageCriterion damage_crit{EPS_YIELD};
DamageTracker<StructModel> tracker{model, damage_crit};
solver.add_observer(tracker);

auto director =
    make_damage_threshold_director<StructModel>(damage_crit, 1.0);
solver.step_to(T_MAX, director);   // pauses at first yield

// --- Transition report ---
auto report = tracker.transition_report();
// report->trigger_time, report->critical_element

// --- Phase 2: extract kinematics, build sub-models, solve ---
for (auto e_idx : top3_damaged_beam_elements) {
    const auto* beam = model.elements()[e_idx].as<BeamElemT>();
    auto u_e   = model.elements()[e_idx].extract_element_dofs(u_local);
    auto kin_A = extract_section_kinematics(*beam, u_e, -1.0);
    auto kin_B = extract_section_kinematics(*beam, u_e, +1.0);
    coordinator.add_critical_element({e_idx, kin_A, kin_B, xA, xB});
}
coordinator.build_sub_models({COL_B, COL_H, 4, 4, 8});

SubModelSolver sub_solver{EC_COL, NU_RC};
for (std::size_t i = 0; i < coordinator.sub_models().size(); ++i) {
    auto result = sub_solver.solve(coordinator.sub_models()[i],
                                   OUT + "sub_models/sub_" +
                                   std::to_string(i));
    auto hs = homogenize(result, coordinator.sub_models()[i],
                         COL_B, COL_H);
    // hs.N, hs.Vy, hs.Mz, hs.E_eff, hs.G_eff
}
\end{lstlisting}

% ---------------------------------------------------------------------------
\subsection{Numerical results}
\label{sec:seismic_results}
% ---------------------------------------------------------------------------

\paragraph{Phase 1 outcome.}
At the time of writing, the dynamic run has \emph{not yet produced a valid}
first-yield time $t_y$.  This is no longer interpreted as a physical
``elastic response'' result; it is now understood as a diagnostic issue in
the dynamic post-processing path.  The chronology is important:

\begin{enumerate}
  \item The first working seismic executable used the MYG012 record, but the
    code never called \texttt{solver.set\_time\_step(DT)} before
    \texttt{step\_to()}.  PETSc therefore advanced with its own default
    time step and heavy step subdivision.  The original argument
    ``$\varepsilon_{\max}=0$ despite large PGA, therefore the section is
    not updating'' was \emph{wrong}; the run had not yet reached the
    strong-motion window at the expected physical time.
  \item After forcing the intended time step, the diagnostics
    \texttt{|u|\_inf} showed nonzero structural motion, so the earlier
    claim ``the structure is not responding'' was also false.
    What was actually happening was severe time-step compression near the
    P-wave arrival.
  \item A second incorrect assumption was that adding
    \texttt{-snes\_max\_it}, \texttt{-snes\_rtol}, \texttt{-snes\_atol},
    and \texttt{-ksp\_max\_it} after \texttt{solver.setup()} would affect
    the already-created PETSc solver objects.  PETSc later reported these
    options as unused.  This was corrected by moving the options before
    \texttt{setup()} and by setting tolerances directly on the live
    \texttt{SNES/KSP} objects.
  \item The excitation was then changed from the late-peaking MYG012 window
    to the stronger near-fault MYG004 record, and the scale factor was
    made a runtime parameter to allow controlled sweeps.
  \item The decisive instrumentation step was to print, at every 500 steps,
    the global displacement norm $|u|_\infty$, the local-state norm
    $|u_\text{local}|_\infty$, the maximum section generalized strain, and
    the maximum fibre strain reconstructed from \texttt{section\_snapshots()}.
    These diagnostics showed a consistent pattern:
    \begin{itemize}
      \item $|u|_\infty > 0$ in the PETSc time-step monitor,
      \item $|u_\text{local}|_\infty = 0$ after localization,
      \item section generalized strains remain zero,
      \item reconstructed fibre strains remain zero.
    \end{itemize}
    At this point the localization mechanism (PETSc section layout,
    constraint registration, \texttt{DMGlobalToLocal} API usage) was
    verified correct in isolation: a manual scatter inside the director
    lambda produced $|u_\text{local}| = 2.9\times10^{-4} > 0$.
    The problem was therefore \emph{not} in the scatter itself.

  \item The root cause was identified by instrumenting
    \texttt{post\_step\_()} with \texttt{printf}+\texttt{fflush} calls
    that \emph{never appeared in output}, even though the observer
    callback invoked \emph{from within} \texttt{post\_step\_()} was
    known to execute (it produced diagnostic output through the director
    lambda).  This proved that the \texttt{TSMonitor} callback
    registered via \texttt{TSMonitorSet()} was \textbf{not being fired}
    by PETSc during \texttt{TSStep()} on this platform
    (Windows/MSYS2, PETSc~3.22).

    The architecture of \texttt{step\_to()} was:
    \begin{enumerate}[label=(\alph*)]
      \item Call \texttt{step()} $\to$ \texttt{TSStep()} advances the
            global solution $\mathbf{U}^{n+1}$.
      \item Call the director callback to evaluate damage and decide
            whether to pause.
    \end{enumerate}
    Step~(b) read \texttt{model\_->state\_vector()}, which is a
    \emph{local} PETSc vector that was supposed to be updated by
    \texttt{post\_step\_()} between (a) and (b).  Since the
    \texttt{TSMonitor} never fired, the local state was never updated,
    and the damage criterion always saw identically~zero.

    \textbf{Fix.}  The global-to-local scatter, imposition of Dirichlet
    values, and material state commitment were moved \emph{directly}
    into \texttt{step\_to()} and \texttt{step\_n()}, immediately after
    \texttt{step()} and before the director call:
    %
\begin{lstlisting}[language=C++,basicstyle=\ttfamily\scriptsize]
// Inside step_to(), after step():
DM dm;
TSGetDM(ts_, &dm);
Vec U_cur;
TS2GetSolution(ts_, &U_cur, nullptr);  // global solution
Vec scratch;
DMGetLocalVector(dm, &scratch);
DMGlobalToLocal(dm, U_cur, INSERT_VALUES, scratch);
VecAXPY(scratch, 1.0, model_->imposed_solution());
VecCopy(scratch, model_->state_vector());
DMRestoreLocalVector(dm, &scratch);
model_->commit_material_state(model_->state_vector());
// NOW call director -- it sees updated local state
\end{lstlisting}
    %
    This makes \texttt{step\_to()} self-contained with respect to state
    synchronisation, regardless of whether \texttt{TSMonitor} fires.
\end{enumerate}

\paragraph{Validation.}
With the fix applied, the dynamic analysis produces physically consistent
results.  Two scale-factor levels were tested:

\begin{itemize}
  \item \textbf{Scale $= 2.0$}: peak damage index $d_{\max} = 0.663$
        (66.3\% of yield strain); no yielding detected in 100\,s.
        This confirms the frame remains elastic under moderate
        amplification.
  \item \textbf{Scale $= 5.0$}: \textbf{first yield detected} at
        $t_y = 65.24\,\text{s}$, with damage index $d = 1.034$
        (fibre strain $\varepsilon = 2.172 \times 10^{-3} >
        \varepsilon_y = 2.1 \times 10^{-3}$).  The critical elements
        are columns~0 and~2 (first-storey columns), consistent with the
        expected soft-storey mechanism under strong horizontal excitation.
\end{itemize}

The trigger time $t_y = 65.24\,\text{s}$ corresponds to $t = 95.24\,\text{s}$
in the original (untrimmed) record, which falls within the peak-PGA zone
of the Tohoku main shock ($\text{PGA} = 26.794\,\text{m/s}^2$ at the
MYG004 station).  The first-storey column failure mode is physically
expected: the fixed-base columns carry the full horizontal shear from
all three storeys while having the same cross-section as the upper
columns.

\paragraph{Phase 2 sub-model results.}

Phase~2 sub-model extraction is currently blocked by a secondary bug:
the section-level generalised strains ($\varepsilon_0$, $\kappa_z$)
read from \texttt{FiberSection3D::current\_state()} return zero despite
the fibre-level strains being nonzero after commit.  This prevents the
beam-theory displacement field from being imposed as boundary conditions
on the continuum sub-models.  The table below is preserved as pending
output; the sub-model kinematics extraction requires further
investigation.

\begin{table}[htb]
\centering
\caption{Pending sub-model results at first valid yielding (column elements,
     elastic sub-model $E_c = 25\,740\,\text{MPa}$,
     $\nu = 0.20$, $4\times4\times8$ mesh).}
\label{tab:seismic_submodel_results}
\begin{tabular}{lrrrrr}
\hline
Sub-model & $N$ [kN] & $V_y$ [kN] & $M_z$ [kN$\cdot$m]
          & $E_\text{eff}$ [GPa] & $G_\text{eff}$ [GPa] \\
\hline
0 (col.) & --- & --- & --- & --- & --- \\
1 (col.) & --- & --- & --- & --- & --- \\
2 (col.) & --- & --- & --- & --- & --- \\
\hline
\end{tabular}
\end{table}

Once the section kinematics extraction is corrected, the effective
Young's modulus $E_\text{eff}$ should recover the ACI initial tangent
stiffness $E_c = 4700\sqrt{f'_c}$ since the continuum sub-model uses
linear elasticity; deviations below $2\%$ are expected from mesh-size
effects as demonstrated in the elastic verification
(Section~\ref{sec:elastic_multiscale_example}).

% ---------------------------------------------------------------------------
\subsection{Infrastructure footprint}
% ---------------------------------------------------------------------------

The seismic example exercises every component added during Phases~0--6:

\begin{itemize}
  \item \textbf{Phase 0}: \texttt{GroundMotionRecord::from\_knet()} parses
        the K-NET text file; \texttt{trim()} extracts the first 120~s.
  \item \textbf{Phase 1}: \texttt{LShapedBuildingDomainBuilder} (here used
        as a rectangular sub-case) constructs nodes and element connectivity.
  \item \textbf{Phase 2}: \texttt{DamageTracker} + \texttt{TransitionDirector}
        drive the linear-to-nonlinear transition without modifying the solver.
  \item \textbf{Phase 3}: \texttt{PrismaticDomainBuilder} generates one
        $4\times4\times8$ hex8 mesh per critical element.
  \item \textbf{Phase 4}: \texttt{FieldTransfer::transfer()} imposes the
        beam-theory displacement field as Dirichlet BCs on the sub-model
        boundary faces.
  \item \textbf{Phase 5}: \texttt{MultiscaleCoordinator} manages element
        registration, BC computation, and sub-model life-cycle.
  \item \textbf{Phase 6}: VTK output via \texttt{SubModelSolver::solve()}
        with \texttt{vtk\_prefix} and global-frame export via
        \texttt{fall\_n::vtk::StructuralVTMExporter}.
\end{itemize}

% ---------------------------------------------------------------------------
\subsection{Multiscale functionality summary}
% ---------------------------------------------------------------------------

Table~\ref{tab:multiscale_summary} lists all classes and free functions
that implement the multiscale framework, together with the phase in which
they were introduced.

\begin{table}[htb]
\centering
\caption{Multiscale framework -- complete functionality inventory.}
\label{tab:multiscale_summary}
\small
\begin{tabular}{lll}
\hline
Class / function & Phase & Role \\
\hline
\texttt{GroundMotionRecord}          & 0 & K-NET/COSMOS parser, trim, PGA \\
\texttt{LShapedBuildingDomainBuilder}& 1 & RC frame domain + connectivity \\
\texttt{DamageTracker<ModelT>}       & 2 & per-step element damage monitoring \\
\texttt{TransitionDirector}          & 2 & pause/stop callback for \texttt{step\_to} \\
\texttt{MaxStrainDamageCriterion}    & 2 & fibre strain $\div$ $\varepsilon_y$ \\
\texttt{PrismaticDomainBuilder}      & 3 & hex8 prismatic mesh on PETSc DMPlex \\
\texttt{FieldTransfer}               & 4 & beam kinematics $\to$ hex BC \\
\texttt{ElementKinematics}           & 4 & section kinematics at both element ends \\
\texttt{SectionKinematics}           & 4 & $\varepsilon_0, \kappa_y, \kappa_z$,
                                          displ., rotation, $E$, $G$, $\nu$ \\
\texttt{MultiscaleCoordinator}       & 5 & sub-model life-cycle manager \\
\texttt{SubModelSolver}              & 5 & elastic continuum solve + homogenize \\
\texttt{homogenize()}                & 5 & $\langle\sigma\rangle \to N,V,M$ \\
\texttt{SubModelSolver::solve(,vtk)} & 6 & VTK export of deformed mesh + GP cloud \\
\texttt{StructuralVTMExporter}       & 6 & global frame VTK at yield \\
\hline
\end{tabular}
\end{table}

